\documentclass[12pt]{article}

\usepackage[utf8]{inputenc}
\usepackage[T1]{fontenc}

\title{La lingüística de la biología}
\author{J. Toledo, S. Castro, A. Jara, D. Rojas}
\date{\today}

\begin{document}

\maketitle

\section{Introducción}
<<<<<<< HEAD
Las interacciones humanas han sido la base fundamental en la que se ha desarrollado la sociedad a lo largo de la historia. La tendencia humana de agruparse se ha visto reflejada desde mucho tiempo atrás, siendo la respuesta evolutiva común a la necesidad de supervivencia. El ser humano no se encuentra en la cima de la cadena trófica por su fuerza o físico, sino por su capacidad intelectual. Por lo que, para lograr una cooperación efectiva, transformó la comunicación simple en lenguajes, que poco a poco se fueron especializando hasta lo que conocemos. Resulta mediocre etiquetar un proceso tan trascendental como solo un producto cultural, por lo que debe entenderse como un mecanismo con un fundamento biológico y cognitivo. La biolingúistica busca analizar la emergencia del lenguaje, no solo por esa presión de comunicación, sino también por principios de eficiencia y desarrollo cognitivo (tercer factor), dando a entender que es una propiedad biológica propia y una capacidad física del cerebro.
Esta perspectiva surgida de los trabajos de Eric Lenneberg y Noam Chomsky, considera que el tercer factor está formado por principios cognitivos mucho más amplios que regulan la arquitectura interna de dicho sistema, mediante criterios de eficiencia computacional (analizando el cerebro como un procesador biológico) y minimo esfuezo (Vitral L. 2007). Donde, ante una situación que ponga a elegir entre minimizar el costo del cálculo cerebral (interno) y facilitar la comunicación (externo), se prioriza el primero.
Pero para lograr el avance colectivo, es necesario compartir la información que se procesa internamente, por lo que debe expresarse de tal manera que pueda ser interpretada sin que pierda sentido. Aquí tiene un papel fundamental la macrosintaxis y la lingúistica pragmática.
Fuentes Rodriguez (2017) argumenta que el lenguaje no se limita a un conjunto de oraciones articuladas alrededor del verbo y sus complementos, sino que transiende a un mensaje real donde el nucleo (o "noyau") y el contexto comunicativo se rigen por criterios entonativos, informativos y pragmaticos que logran una correcta interpretación del mensaje.

La relación entre biología y lingüística aplicada se evidencia en la manera en que los procesos comunicativos humanos no solo responden a necesidades sociales, sino también a condicionantes practicas, como en el caso a observar que es la lingüística de la biología, este tipo de lingüística fue necesaria para abarcar un problemas, el lenguaje que se aplicaría a  conceptos, textos, descubrimientos, etc.    

=======
Las interacciones humanas han sido la base fundamental sobre la que se ha desarrollado la sociedad a lo largo de la historia. La tendencia humana de agruparse se ha visto reflejada desde mucho tiempos remotos, siendo una respuesta evolutiva común a la necesidad de supervivencia. El ser humano no se encuentra en la cima de la cadena trófica por su fuerza o físico, sino por su capacidad intelectual. Por ello, para lograr una cooperación efectiva, transformó la comunicación simple en lenguajes, que poco a poco se fueron especializando hasta llegar a lo que conocemos actualmente. Resulta mediocre etiquetar un proceso tan trascendental como un simple producto cultural, por lo que debe entenderse como un mecanismo con un fundamento biológico y cognitivo. 
La biolingúística busca analizar la emergencia del lenguaje, no solo a partir de la presión comunicativa, sino también por principios de eficiencia y desarrollo cognitivo (tercer factor), dando a entender que constituye una propiedad biológica propia y una capacidad física del cerebro.
Esta perspectiva surgida de los trabajos de Eric Lenneberg y Noam Chomsky, considera que el tercer factor está formado por principios cognitivos mucho más amplios que regulan la arquitectura interna de dicho sistema, mediante criterios de eficiencia computacional (analizando el cerebro como un procesador biológico) y mínimo esfuezo (Vitral L. 2007). De este modo ante una situación que implique elegir entre minimizar el costo del cálculo cerebral (interno) y facilitar la comunicación (externo), se prioriza el primero.
Pero, para lograr el avance colectivo, es necesario compartir la información que se procesa internamente, por lo que debe expresarse de tal manera que pueda ser interpretada sin que pierda sentido. Aquí tienen un papel fundamental la macrosintaxis y la lingüística pragmática.
La relación entre biología y lingüística aplicada se evidencia en la manera en que los procesos comunicativos humanos no solo responden a necesidades sociales, sino también a condicionantes prácticos. Este es el caso de la lingüística de la biología; que fue necesaria para abarcar un problema: el lenguaje que se aplicaría a conceptos, textos y descubrimientos, etc.    
Como la biología es una ciencia basada en el desarrollo de conocimientos, resulta crucial esta relación para abordar la forma en que se construye el lenguaje científico; siendo necesario un  manejo adecuado de las reglas macrosintácticas y pragmáticas que organizan el mensaje en el contexto, y a su vez, el concientizar las restricciones biologicas y cognitivas en el procesamiento interno de la información.  
>>>>>>> 66ac432b257765006ed4e51e45b0a92fadc0e4e5
Referencias
Vitral, L.(2017). La emergencia de la biolingüística y sus consecuencias para la teoría lingúística. Lingüística y literatura.
\end{document}
